\documentclass[a4paper,fontset=windows]{ctexart}

\usepackage{anyfontsize}
\usepackage{amsmath,amssymb,mathrsfs}
\usepackage{autobreak}
\allowdisplaybreaks
\usepackage{indentfirst}
\setlength{\parindent}{0em}
\usepackage{geometry}
\geometry{top=3cm,bottom=3cm,left=2cm,right=2cm}
\usepackage{setspace}
\usepackage{bm}
\usepackage{graphicx,caption,subcaption}
\renewcommand{\arraystretch}{1.25}


\begin{document}

\title{\textbf{实验三十\ \ 用示波器观测动态磁滞回线}}
\author{1900011604 张植竣}
\date{2021年5月31日}

\maketitle

\section*{1\ \ 实验数据和现象}

\subsection*{1.1\ \ $100\,\mathrm{Hz}$下铁氧体的饱和磁滞回线}

\begin{table}[!ht]
    \centering
        \begin{tabular}{|c|c|c|c|c|c|c|c|c|c|}
            \hline
            $u_x/\mathrm{mV}$ & $-215$ & $-200$ & $-175$ & $-150$ & $-125$ & $-100$ & $-75$ & $-50$ & $-25$ \\
            \hline
            $u_y/\mathrm{mV}$ & $-13.45$ & $-13.40$ & $-13.15$ & $-12.85$ & $-12.50$ & $-11.95$ & $-11.15$ & $-10.00$ & $-7.45$ \\
            \hline
            $u_x/\mathrm{mV}$ & $-215$ & $-200$ & $-175$ & $-150$ & $-125$ & $-100$ & $-75$ & $-50$ & $-25$ \\
            \hline
            $u_y/\mathrm{mV}$ & $-13.45$ & $-13.35$ & $-12.95$ & $-12.55$ & $-11.75$ & $-10.65$ & $-8.65$ & $-5.60$ & $-1.10$ \\
            \hline
            $u_x/\mathrm{mV}$ & 25 & 50 & 75 & 100 & 125 & 150 & 175 & 200 & 215 \\
            \hline
            $u_y/\mathrm{mV}$ & 7.70 & 10.30 & 11.60 & 12.55 & 13.10 & 13.55 & 13.80 & 14.10 & 14.15 \\
            \hline
            $u_x/\mathrm{mV}$ & 25 & 50 & 75 & 100 & 125 & 150 & 175 & 200 & 215 \\
            \hline
            $u_y/\mathrm{mV}$ & 1.35 & 5.55 & 9.00 & 11.10 & 12.35 & 13.20 & 13.65 & 14.05 & 14.15 \\
            \hline
        \end{tabular}
\end{table}

\begin{figure}[!ht]
    \centering
    \includegraphics[width=0.8\linewidth]{1.png}
    \caption{$100\,\mathrm{Hz}$下铁氧体的饱和磁滞回线}
\end{figure}

\begin{table}[!ht]
    \centering
    \begin{tabular}{|c|c|c|c|c|c|}
        \hline
        $\left.2u_y\right|_{B_s}/\mathrm{mV}$ & $\left.2u_y\right|_{B_r}/\mathrm{mV}$ & $\left.2u_x\right|_{H_c}/\mathrm{mV}$ & $B_s/\mathrm{T}$ & $B_r/\mathrm{T}$ & $H_c/\mathrm{A \cdot m^{-1}}$ \\
        \hline
        27.60 & 7.20 & 37.0 & 0.3710 & 0.0968 & 10.67 \\
        \hline
    \end{tabular}
\end{table}

\subsection*{1.2\ \ 不同频率下铁氧体的饱和磁滞回线}

原始数据如下：

\begin{table}[!ht]
    \centering
    \begin{tabular}{|c|c|c|c|c|}
        \hline
        $f/\mathrm{Hz}$ & $\left.2u_y\right|_{B_r}/\mathrm{mV}$ & $\left.2u_x\right|_{H_c}/\mathrm{mV}$ & $B_r/\mathrm{T}$ & $H_c/\mathrm{A \cdot m^{-1}}$ \\
        \hline
        50 & 7.16 & 38.1 & 0.0962 & 10.99 \\
        \hline
        100 & 7.30 & 38.8 & 0.0981 & 11.19 \\
        \hline
        150 & 7.36 & 38.9 & 0.0989 & 11.22 \\
        \hline
    \end{tabular}
\end{table}

误差分析：

实测线宽的一半约$0.5\,\mathrm{mm}$，而示波器一大格是$10\,\mathrm{mm/DIV}$，$u_x$和$u_y$的偏转因数分别为$K_x=10\,\mathrm{mV/DIV}$，$K_y=2\,\mathrm{mV/DIV}$。则线宽导致的误差为：
\[
e_x = \frac{0.5\,\mathrm{mm}}{10\,\mathrm{mm/DIV}}\times10\,\mathrm{mV/DIV} = 0.5\,\mathrm{mV},\quad
\sigma_x = \frac{e_x}{\sqrt{3}} = 0.288675135\,\mathrm{mV}
\]
\[
e_y = \frac{0.5\,\mathrm{mm}}{2\,\mathrm{mm/DIV}}\times10\,\mathrm{mV/DIV} = 0.1\,\mathrm{mV},\quad
\sigma_y = \frac{e_y}{\sqrt{3}} = 0.057735027\,\mathrm{mV}
\]
又因为示波器测量的不连续性，极限误差估计为光标移动一步的大小，即分度值的0.01倍：
\[
e'_x = 0.1\,\mathrm{mV},\quad
e'_y = 0.02\,\mathrm{mV},\quad
\]
考虑到示波器的测量精度（允差）为测出值的$2\%$加满刻度值的$0.3\%$，由$100\,\mathrm{Hz}$的数据可以估计$B_r$和$H_c$的精度为：（因为有两个光标，$e'_x$和$e'_y$需要算两次）
\[
\sigma_{B_r} = 0.001790479,\quad
\sigma_{H_c} = 0.199001948,\quad
\]
测量得到的数据为：

\begin{table}[!ht]
    \centering
    \begin{tabular}{|c|c|c|c|c|}
        \hline
        $f/\mathrm{Hz}$ & $\left.2u_y\right|_{B_r}/\mathrm{mV}$ & $\left.2u_x\right|_{H_c}/\mathrm{mV}$ & $B_r/\mathrm{T}$ & $H_c/\mathrm{A \cdot m^{-1}}$ \\
        \hline
        50 & 7.16 & 38.1 & $0.0962\pm0.0018$ & $10.99\pm0.20$ \\
        \hline
        100 & 7.30 & 38.8 & $0.0981\pm0.0018$ & $11.19\pm0.20$ \\
        \hline
        150 & 7.36 & 38.9 & $0.0989\pm0.0018$ & $11.22\pm0.20$ \\
        \hline
    \end{tabular}
\end{table}

磁滞回线的面积对应于循环磁化一周所发生的能量损耗。对材料进行准静态磁化时，损耗来自磁滞损耗；对材料进行交流动态磁化时，除了有磁滞损耗外，还会有涡流损耗和剩余损耗。频率增高时损耗增大，而饱和磁场强度、饱和磁感应强度不变（这是材料的内禀属性），磁滞回线面积增大要求$B_r$和$H_c$增大。金属氧化物组成的铁氧体材料电阻率高，高频条件下其涡流损耗很小，则$B_r$和$H_c$的值随频率升高基本不变，仅略微增大。

从测量数据中可以看出$B_r$和$H_c$的值随频率变化的规律符合理论推测。

\subsection*{1.3\ \ 不同积分常量下的Lissajous图形}

\begin{figure}[ht!]
    \centering
    \begin{subfigure}[t]{0.3\linewidth}
        \centering
        \includegraphics[width=1\linewidth]{RC=0.01s.jpg}
        \subcaption*{$R_2=2\,\mathrm{k}\Omega,\ C=5\,\mu\mathrm{F},\ \tau=0.01\,\mathrm{s}$}
    \end{subfigure}
    \quad
    \begin{subfigure}[t]{0.3\linewidth}
        \centering
        \includegraphics[width=1\linewidth]{RC=0.05s.jpg}
        \subcaption*{$R_2=10\,\mathrm{k}\Omega,\ C=5\,\mu\mathrm{F},\ \tau=0.05\,\mathrm{s}$}
    \end{subfigure}
    \quad
    \begin{subfigure}[t]{0.3\linewidth}
        \centering
        \includegraphics[width=1\linewidth]{RC=0.5s.jpg}
        \subcaption*{$R_2=100\,\mathrm{k}\Omega,\ C=5\,\mu\mathrm{F},\ \tau=0.5\,\mathrm{s}$}
    \end{subfigure}
\end{figure}

积分常量$\tau=R_2 C$影响电容$C$上的电压$\displaystyle{u_y=u_C=\frac{1}{\tau}\int{u_{R_2}}\,\mathrm{d}t}$，若不满足$\tau \gg T$的条件，那么$u_{R_2}=u_2$就无法成立，且$u_y$的相位改变导致Lissajous图形扭曲，无法反映真实的磁滞回线。但是积分常量不影响真实磁滞回线的形状，这个是由材料本身决定，与外界测量电路条件无关的。

\subsection*{1.4\ \ $100\,\mathrm{Hz}$下铁氧体的动态磁化曲线}

\begin{table}[!ht]
    \centering
    \begin{tabular}{|c|c|c|c|c|c|c|c|c|c|}
        \hline
        $u_x/\mathrm{mV}$ & 0.00 & 1.00 & 2.00 & 3.00 & 4.00 & 5.00 & 10.00 & 15.0 & 20.0 \\
        \hline
        $u_y/\mathrm{mV}$ & 0.000 & 0.080 & 0.160 & 0.240 & 0.320 & 0.400 & 0.880 & 1.40 & 1.96 \\
        \hline
        $u_x/\mathrm{mV}$ & 25.0 & 30.0 & 35.0 & 40.0 & 50.0 & 60.0 & 70.0 & 80.0 & 90.0 \\
        \hline
        $u_y/\mathrm{mV}$ & 2.58 & 3.22 & 3.86 & 4.58 & 5.85 & 7.15 & 8.30 & 9.10 & 10.05 \\
        \hline
        $u_x/\mathrm{mV}$ & 100.0 & 150.0 & 200.0 & 300.0 & 400.0 & 500.0 & 600.0 & 700.0 & 800.0 \\
        \hline
        $u_y/\mathrm{mV}$ & 10.85 & 12.85 & 13.70 & 14.50 & 14.85 & 15.15 & 15.35 & 15.50 & 15.65 \\
        \hline
    \end{tabular}
\end{table}

\begin{figure}[!ht]
    \centering
    \includegraphics[width=0.8\linewidth]{2.png}
    \caption{$100\,\mathrm{Hz}$下铁氧体的动态磁化曲线}
\end{figure}

\begin{figure}[!ht]
    \centering
    \includegraphics[width=0.8\linewidth]{3.png}
    \caption{$100\,\mathrm{Hz}$下铁氧体的$\mu_m-H_m$曲线}
\end{figure}

起始磁导率为$\mu_i=2966.328688451$，$H=37.0078\,\mathrm{A\cdot m^{-1}}$处$\mu_m$达到最大，为$4432.7478$。

\subsection*{1.5\ \ 不同频率下硅钢样品在给定交变磁场幅度$H_m=400\mathrm{A\cdot m^{-1}}$下动态磁滞回线如何变化}

\begin{table}[!ht]
    \centering
    \begin{tabular}{|c|c|c|c|c|c|c|}
        \hline
        $f/\mathrm{Hz}$ & $\left.2u_y\right|_{B_r}/\mathrm{mV}$ & $\left.2u_y\right|_{B_m}/\mathrm{mV}$ & $\left.2u_x\right|_{H_c}/\mathrm{mV}$ & $B_r/\mathrm{T}$ & $B_m/\mathrm{T}$ & $H_c/\mathrm{A \cdot m^{-1}}$ \\
        \hline
        20 & 42.4 & 68.1 & 209.0 & 0.5889 & 0.9458 & 104.50 \\
        \hline
        40 & 44.6 & 67.7 & 244.5 & 0.6194 & 0.9403 & 122.25 \\
        \hline
        20 & 46.1 & 67.7 & 282.0 & 0.6403 & 0.9403 & 141.00 \\
        \hline
    \end{tabular}
\end{table}

磁滞回线的面积对应于循环磁化一周所发生的能量损耗。对材料进行准静态磁化时，损耗来自磁滞损耗；对材料进行交流动态磁化时，除了有磁滞损耗外，还会有涡流损耗和剩余损耗。频率增高时损耗增大，而饱和磁感应强度$B_m$不变（这是材料的内禀属性），磁滞回线面积增大要求$B_r$和$H_c$增大。合金组成的硅钢电阻率低，高频磁化时其涡流损耗大，则$B_r$和$H_c$的值随频率升高而明显增大。

从测量数据中可以看出$B_m$基本不变，$B_r$和$H_c$的值随频率升高而增大，符合理论推测。

\section*{2\ \ 思考题}

1. 静态磁滞回线：准静态磁化铁磁体时铁磁体的磁滞回线，即外加磁场$H$变化的速率很小，样品有充分的时间响应，没有涡旋电动势产生；动态磁滞回线：动态态磁化铁磁体时铁磁体的磁滞回线，即外加磁场$H$的变化有一定的速度，样品每一时刻的磁化并没有达到稳定，有涡旋电动势产生。铁磁材料动态磁滞回线的形状和面积与材料本身的性质、外加磁场的频率和振幅有关。

2. 磁滞回线的面积对应于循环磁化一周所发生的能量损耗。对材料进行准静态磁化时，损耗来自磁滞损耗；对材料进行交流动态磁化时，除了有磁滞损耗外，还会有涡流损耗和剩余损耗。频率增高时损耗增大，而饱和磁场强度、饱和磁感应强度不变（这是材料的内禀属性），磁滞回线面积增大要求$B_r$和$H_c$增大。金属氧化物组成的铁氧体材料电阻率高，高频条件下其涡流损耗很小，则$B_r$和$H_c$的值随频率升高基本不变，仅略微增大。合金组成的硅钢材料电阻率低，高频磁化时其涡流损耗大，则$B_r$和$H_c$的值随频率升高而明显增大。

3. 需满足$\tau=R_2 C \gg T$的条件，其中$T=f^{-1}$为磁场（励磁电流）的周期。

4. 将磁场频率调低以至于肉眼无法看到完整的磁滞回线而是一个不断绕行的磁滞回线，即可确定磁滞回线的绕行方向。用高速快门（快门时间$\Delta t < T$）拍摄磁滞回线也可以达到相同效果。

\end{document}